% !TeX root = ../papers/nc-p3.tex
\subsection{正次元の点対応でも検出できるOre窓と楕円基底の分類}
\label{sec:elliptic-ore-recognition}
\begin{proposition}\label{prop:ore-mark-recognition}
標準有限Frobenius四窓の中で、定理\ref{thm:maximal-tangent-ore}の三生成Ore表示を持つ
標準AS対象とその指定点を、辺の有限性を仮定せずに検出できる。
標準四窓上で四点 $[p_i]\in\mathbb P(V_i^*)$ を取り、
各 $(p_i,p_{i+1})$ が辺の零点で、関係の左右評価が個別に階数3、
二評価を合わせた射影Jacobi行列の階数も3という条件を課す（添字は法4）。
この条件の下で四生成候補の標準AS正則性は、復元した三生成基底の標準二次AS正則性に同値である。
点を指定した同型と移動は $P_B$ と $F_B$ で決まる。
\end{proposition}
\begin{proof}
点を反復すれば全整数上で同じ評価条件を得る。定理\ref{thm:maximal-tangent-ore}で
$S_i$ と可逆な $f_i$ を得る部分はこのJacobi行列だけを使う。
純粋な三次交叉の次元が1となる証明で用いたのも、四生成三次元数20と
直前の点での評価階数3だけである。その交叉の両端支持も連続する指定点での階数3である。
従って高次の直和表示・正則性の両方向・指定同型の証明が全て適用される。
逆に標準ASなOre表示では $H_i$ の零化点の両評価とJacobi行列が上記階数を持つ。
この列を標準四窓へ運ぶと、同定前の列に対する純粋交叉の末端支持の証明から
$[p_{i+2}]=[p_{i-2}]$ が従うので、実際に四点の検査で拾える。
有限性から得た「高々二本」という個数をここでは使っていない。
\end{proof}

\begin{theorem}\label{thm:all-elliptic-ore-windows}
上の検出条件を満たし、三生成部分の三関係が滑らかな平面三次曲線 $E$ の
自己同型 $\sigma$ のグラフを定める標準AS対象を考える。正次元の辺も許す。
指定点は対象から一意に復元できる。$L=\mathcal O_E(1)$ として基底は
\[
        L\ne\sigma^* L,\qquad L\otimes(\sigma^2)^*L\simeq(\sigma^*L)^2
\]
を満たす $(E,L,\sigma)$ で尽くされる。各基底に対する補正は次の通りである。
\begin{enumerate}
\item $\sigma$ が位数2,4の平行移動なら定理\ref{thm:elliptic-tangent-classification}の
九次元空間である。他の平行移動の補正は零だけである。
\item 固定点があるとき、原点保持部分 $\alpha$ の位数は2,3,4のいずれかである。
位数2では補正空間は二次元、位数3,4では零である。
\end{enumerate}
零補正は常に $\check{B[Z]}$ である。有限四辺の部分以外は全て最小周期1である。
添字保持同型類を得るには基底を偏極曲線と移動の同型で比較し、
非零補正はその射影類を有限群 $\Aut(E,L,\sigma)$ の誘導作用で比較すればよい。
\end{theorem}
\begin{proof}[基底と指定点]
二つの束条件の必要性は定理\ref{thm:elliptic-tangent-classification}と
補題\ref{lem:reduced-ternary-fourth-shift}の行列式計算による。
逆にこのデータから $S=\ker(H^0(L)\otimes H^0(L)\to H^0(L\otimes\sigma^*L))$ を取る。
第二因子の切断は $\sigma$ で引き戻す。前節のpencilの完全列から $\dim S=3$ である。
平面の点 $p$ の課す二切断の鉛筆は、$p\notin E$ なら $E$ 上で基点自由であり、
$\sigma^*L\otimes L^{-1}$ に切断がないことから $S$ の $p$ での評価核が零になる。
$p\in E$ なら $L(-p)$ の基点自由な二切断を使い、
$\sigma^*L\otimes L^{-1}(p)$ の切断が一次元なので評価階数は2である。
従って行列式は $E$ の三次方程式であり、対応はスキームとして所定のグラフとなる。

定理\ref{thm:elliptic-tangent-classification}の $n=3$ の束の列で、
中央写像の核直線束は行列式から
$L\otimes(\sigma^2)^*L\otimes(\sigma^*L)^{-2}\simeq\mathcal O_E$ となる。
その生成切断は $S\otimes H$ に入り、像が零なので $H\otimes S$ にも入る。
従って一次元の三次交叉を生成する。片側縮約が全三次元を張ることも確認する。
$p$ に消える核関係の直線 $E\to\mathbb P(S)$ は単射である。
二点で同じ関係の左評価が零なら、その行列は階数1以下である。
しかし可分解な非零関係が既約曲線のグラフ上で消えることはない。
従ってこの写像の像は直線に入らない（楕円曲線から射影直線への単射はない）。
逆側も同じである。これで先の束の完全列からbarの初項、三次の非零因子、
Hilbert級数とKoszul--Frobenius性を得た証明がそのまま適用され、$B$ は標準三次元ASである。

指定点以外の辺零点では両端で三生成部分への射影が定まる。
片方だけが零となることは混合関係が禁じる。
局所的に基底の両射影の自明化を選ぶと $(p,\sigma p)$ の上で増えるのは第四座標二つであり、
混合関係の一行で一つを消せる。従って局所的に $E\times\mathbb A^1$ の部分スキームに入り、
接空間は高々二次元である。三次元の指定点は唯一で、正次元の辺でも指定同型条件が使える。
\end{proof}

\begin{proof}[移動と有理切断による補正]
$\sigma$ が固定点を持たない場合は平行移動である。$\sigma^4\ne1$ のとき
Hesse表示で三次交叉は巡回対称なので $\nu=1$ である。
前節までの二平面の議論を逆に読めば、任意の可延長補正について
$\det[p,\sigma p,d_i(p,\sigma p)]$ は一般点で零である。
よって全点で零となり、定理\ref{thm:elliptic-tangent-classification}の束の完全列から
初期補正は左右乗法で尽きる。位数2,4の結果は同定理による。

固定点があるときも、$\nu=\sigma^4$ が全曲線上で成立すれば固定点で
行列式の微分が消えて滑らかさに反する。従ってやはり上の行列式は恒等零である。
$q=\sigma p$ とし、一般点での二ベクトルの独立性により
\[
    d_i(h)(p,q)=a_i(p)h(q)-h(p)b_{i+1}(q)
\]
と一意に書ける。$a_i,b_{i+1}$ は $L$ の有理切断である。
固定点以外では二行を解けるので正則である。
$\sigma$ の不変微分に対する倍率を $\eta\ne1$ とすれば、各固定点の微分も同じ $\eta$ である。
その近傍では $p$ と $q$ の二次小行列のうち一つは単純零を持つ。
実際、埋込みの非零接ベクトルに対し二つの微分の差は $1-\eta$ 倍である。
従って極は高々単純である。
三次の混合条件を $(p,q,\sigma q)$ で評価すると内側の係数の差
$a_{i+1}(q)-b_{i+1}(q)$ に各 $s(p,\sigma q)$ が掛かる。
一般点ではこれらの全てが零にはならないので $a_{i+1}=b_{i+1}$ である。

固定点の全てで $a_i$ の主要部を取り、その列を $\xi_i$ とする。
局所座標で $h(p)$ と $h(q)$ は同じ定数項を持つため、
$d_i(h)$ の正則性は $\xi_{i+1}=\eta\xi_i$ を与える。
固定点因子を $T$ とすれば $H^1(L)=0$ により
$\xi_0$ は $H^0(L(T))$ の一元 $a$ の主要部となる。
$a_i-\eta^i a$ は正則で、これらの差は補生成元を変更して消せる。
従って全ての類は $d_i=\eta^i d$ の形に表せる。正確には
\[
        P_B=\ker(\lambda_B+\eta\rho_B)/
                 \{L_v-\eta R_v:v\in H\},\qquad F_B=\eta\,1 .
\]
分子の任意の元は実際に全整数列に延びる。
分母の次元は3であり、左右乗法の部分空間の交叉は零である。
実際、$vh=hu$ を全 $h$ についてグラフで評価すると、一般に $p,q$ は独立なので
$v(p)=u(q)=0$、従って $u=v=0$ となる。定理\ref{thm:ternary-ore-linear-space}の
両核の計算で、純粋に補生成元変更から来るこの固有値の元は上記の分母で尽きる。
$X_i$ を $\eta^{-i}$ 倍すれば混合関係は一定になり、最小周期1を得る。
また行列式が恒等零なので一般の $p$ が持ち上がり、辺は正次元である。
\end{proof}

\begin{proof}[固定点を持つ三つの場合の全計算]
変曲原点で $\sigma=t_P\alpha$ と書く。
$L=\mathcal O(3O)$ なので二つの束条件は $3P\ne0$ と $(1-\alpha)3P=0$ である。
原点保持自己同型の位数は $2,3,4,6$ であり、位数6では $1-\alpha$ の核が自明なので除かれる。
$b$ を三等分点として原点を移す射影変換で $P$ を $(1-\alpha)b$ だけ変えられる。
$\alpha$ の位数2,4の場合、$1-\alpha$ は三等分点上可逆であり、
その核上では3倍が可逆である。従って $P$ を $\alpha$ が固定する非零二等分点に取れる。

$\alpha=-1$ ではHesse座標の前二文字の交換 $G$ が逆元を与え、
$P=[1:1:c]$、$c\ne0$, $c^3\notin\{1,-8\}$ を取れる。
従って関係は前の対称三関係に $1\otimes G$ を施したものである。
前節の三次商の零化行列を使って $\lambda_B-\rho_B$ を計算すると、
$\mathbb Q(c)$ 上の階数は13で、分母のない行の十三次小行列式が
\[
                       2 c^{60}(c-1)^{28}(c^2+c+1)^{28}
\]
となる。全ての許容値で非零なので核次元は5、商は二次元である。

位数4なら $y^2=x^3-x$ で $\alpha(x,y)=(-x,\sqrt{-1}\,y)$ またはその逆とする。
唯一の非零固定二等分点 $(0,0)$ の平行移動を合成した像の斉次座標は
$[x^2:x:\sqrt{-1}\,y]$ となる。元の $[1:x:y]$ との九つの積から関係を取り、
$\lambda_B+\sqrt{-1}\rho_B$ の階数は15で、共役も同じである。

位数3では $y^2=x^3+1$、$\alpha(x,y)=(\omega x,y)$、
$\omega^2+\omega+1=0$ と取る。$3P=(0,1)$ または $(0,-1)$ の全18点を考える。
良い素数109ではこれら18点は相異なる有理点へ還元され、$\omega$ は45または63となる。
検算はこの曲線の全108点でグラフを評価して三関係を復元し、
36通り全てで $\lambda_B+\omega\rho_B$ の階数15を得る。
標数零への持上げを明示する。三等分写像は109上エタールであるため
18点は全てこの計算で尽くされる。評価に用いた曲線の各点も良い還元の滑らかさにより持ち上がる。
九列の評価で階数6の単元小行列を取れるので、標数零の三関係の自由基底が持ち上がる。
三次商の余次元は両標数で10である（剰余では直接行列階数を確認）。
従って二次・三次商の基底を局所的に選んで同じ非零小行列式を持ち上げられる。
階数上界15は三次元の補生成元変更から従うので、標数零で商は零である。
これら全ての厳密計算は \path{computations/verify_elliptic_fixed_ore.py} にある。
非零の射影補正への作用は $d(h)\mapsto Gd(G^{-1}h)$ であり、
基底の関係はグラフから一意なので、所定の有限群の比較が零の場合も含め必要十分である。
\end{proof}
